\section{Discussion}
\label{sec:discussion}

\paragraph{Practical deployment.} The hybrid fusion strategy (Stage~3) provides the most favourable accuracy--cost tradeoff: a single conditioned inference pass (2$\times$ compute) yields $\sim$21\% rotation improvement while preserving translation. The full MC iterative method (Stage~4, 13$\times$ compute) offers diminishing marginal gains and is appropriate primarily for offline applications where sub-3\degree{} rotation accuracy justifies the additional cost. In all cases, hybrid decoupling is essential---na\"ive iteration without it leads to progressive translation degradation regardless of the number of steps.

\paragraph{Generality of the framework.} While our experiments use Pi3X, the decoupled hybrid refinement framework applies in principle to any learned geometry model that (i) accepts pose conditioning as input, and (ii) outputs poses defined up to a gauge ambiguity. These properties are shared by several recent architectures (Table~\ref{tab:benchmark}). The contraction analysis (Section~\ref{sec:theory}) provides a diagnostic: measuring input--output error reduction under corrupted conditioning reveals whether a given model exhibits the contraction property and, if so, at what rate. Models with $\lambda < 1$ are amenable to self-refinement; those with $\lambda \geq 1$ are not.

\paragraph{Limitations.} Several constraints bound the applicability of our method. The accuracy floor ($\sim$3\degree{} ARE) is determined by the contraction rate and the model's irreducible error---it cannot be lowered without stronger external priors or architectural modifications to the conditioning mechanism. Translation accuracy is bounded by the quality of unconditioned inference, since hybrid fusion anchors translations to this source. The convergence is empirical: while iterations 1--3 consistently improve, iteration 4 and beyond can produce orientation flips (Table~\ref{tab:convergence}), necessitating a stopping criterion. Finally, our evaluation is conducted on a single multi-camera installation; generalisation to other wide-baseline deployments with different camera counts and geometries remains to be validated.

\paragraph{Future work.} Several directions merit investigation. First, Pi3X also accepts depth conditioning---monocular depth priors (e.g., Depth Anything~v2) could provide complementary geometric signal without requiring ground truth poses, potentially breaking the rotation floor. Second, partial conditioning (conditioning only high-confidence cameras while leaving others unconditioned) may exploit the model's cross-view attention to propagate reliable pose information to uncertain views. Third, the fixed perturbation magnitude ($\sigma_R = 2\degree{}$) could be replaced by an adaptive noise schedule that decreases across iterations, following the annealing principle from diffusion-based methods. Fourth, sparse geometric anchors (a pre-calibrated camera pair or LiDAR registration points) could provide a stronger initialisation that places the estimate deeper within the contraction basin, accelerating convergence and lowering the fixed point.
